\section{结论}
\label{sec:conclusion}

同一个有界阶乘揭示了编译系统的连续性，也揭示了各阶段不可混同的职责。预处理器展开宏却不验证外部实现；前端把文本变成带绑定与类型的树；LLVM IR 把语句重述为基本块、效果和 SSA 数据依赖；后端在 ISA 与 ABI 约束下选择 RV64 操作和寄存器；汇编器把文本编码为字节，并以符号和重定位保存尚未知的地址；链接器兑现这些义务、安排布局和入口；加载器依据段而非节建立进程映像。

\begin{table}[tb]
\caption{端到端表示账本。}
\label{tab:ledger}
\centering\scriptsize
\begin{tabularx}{\linewidth}{@{}lXXX@{}}
\toprule
表示 & 主结构/新增事实 & 淡化或丢弃 & 留给下一层\\
\midrule
源码/预处理文本 & 名称、语句、展开记号 & 注释、宏拼写 & 文法与类型\\
带类型 AST & 树、绑定、类型、值类别 & 排版 & 控制/数据实现\\
LLVM IR & CFG、效果、SSA 定义--使用 & \code{while} 与多数源码形状 & 寄存器、指令、地址\\
RV64 汇编 & 目标操作、ABI 动作、标签 & SSA 名与多数类型 & 编码与符号地址\\
可重定位 ELF & 节、符号、重定位、字节 & 指令文本外观 & 全局解析与布局\\
可执行 ELF & 入口、程序头、已解析布局 & 多数链接义务 & 映射与初始状态\\
进程映像 & 地址、权限、机器状态 & 多数节名和符号名 & 实际执行\\
\bottomrule
\end{tabularx}
\end{table}

表~\ref{tab:ledger} 也修正“降低只会丢信息”的直觉：每层都删除对下一任务无用的区别，同时创造新的、可检查的事实。最终进程里没有宏、WhileStmt 或 \code{\%result}；但它们要求的行为存活在有类型的运算、数据依赖、目标约定、重定位后的位字段和加载映射中。六个测试支持这条具体链的行为一致，计数则只帮助看见结构变化。真正“理解编译系统”，不是背诵命令，而是能在每个边界指出：现在的表示知道什么，为什么知道，以及下一消费者仍需要解决什么。
